\section{引言}

格点上部分子分布函数（PDF）的计算目前是备受关注并投入大量努力的课题
（近期综述及大量文献见 Ref.~\cite{Cichy:2018mum}）。
现代努力的目标是直接提取 PDF $f(x)$ 本身，而非其 $x^N$ 矩。
在格点上，这可以通过从定域算符转向流--流关联函数来实现 \cite{Detmold:2005gg}。
更进一步的想法是从等时关联函数出发 \mbox{\cite{Braun:2007wv}}。

X. Ji 在论文 \cite{Ji:2013dva} 中强烈推动了后续的发展，
提出了一个开创性建议：考虑定义 PDF、分布振幅、广义部分子分布以及横动量依赖分布的
非定域算符的等时版本。
就通常 PDF 而言，Ji 方案的基本概念是``部分子准分布''（quasi-PDF）$Q(y,p_3)$
\cite{Ji:2013dva,Ji:2014gla}，
而 PDF 由准分布在 $p_3 \to \infty$ 大动量极限下得到。

其他方法，如``好格点截面''\cite{Ma:2014jla,Ma:2017pxb}、等时关联函数的 Ioffe 时间分析
\cite{Braun:2007wv,Bali:2017gfr,Bali:2018spj} 以及伪 PDF 方法
\cite{Radyushkin:2017cyf,Radyushkin:2017sfi,Orginos:2017kos}
都以坐标空间为取向，通过取短距离极限 $z_3 \to 0$ 来提取部分子分布。

$p_3\to \infty$ 与 $z_3\to 0$ 这两个极限都是奇异的，
需要使用{\it 匹配关系}
把欧氏格点数据转换为
通常的光锥 PDF。在准 PDF 方法中，此类关系已对夸克
\mbox{\cite{Ji:2013dva,Xiong:2013bka,Ji:2015jwa,Izubuchi:2018srq}} 与胶子
PDF \cite{Wang:2017eel,Wang:2017qyg,Wang:2019tgg}、$\pi$ 介子分布振幅（DA）
\cite{Ji:2015qla} 以及广义部分子分布（GPDs）
\cite{Ji:2015qla,Xiong:2015nua,Liu:2019urm} 作过研究。

在伪 PDF 方法框架内，
非单态 PDF 的匹配关系已在文献
\cite{Ji:2017rah,Radyushkin:2017lvu,Radyushkin:2018cvn,Zhang:2018ggy,Izubuchi:2018srq} 中导出。
利用伪 PDF 方法在格点上提取非单态 GPD 与 $\pi$ 介子 DA 的策略在近期论文
Ref.~\cite{Radyushkin:2019owq} 中作了概述，
其中也导出了这些情形的匹配条件。
本文的主要目标是描述提取非极化胶子 PDF 的伪 PDF 方法的基本要点，并同时求出
单圈匹配条件。

在胶子情形，计算因规范不变性的严格要求而变得复杂。此时，
文献 \cite{Balitsky:1987bk} 的坐标表象方法提供了非常有效的手段。
它基于背景场方法
与热核展开，使我们能够从最初的规范不变双定域算符出发求得单圈修正对其的修改。
结果以显式规范不变的形式得到。

在这一方法中，无需指明特定部分子分布所对应的矩阵元的性质。
这意味着同一套费曼图计算既可用于求
PDF（由前向矩阵元给出）的匹配条件，
也可用于与非前向矩阵元对应的分布振幅和 GPD
（参见 Ref.~\cite{Radyushkin:2019owq} 中对夸克算符如何具体运作的演示）。

本文结构如下。
在第 2 节中，我们分析胶子双定域算符矩阵元的运动学结构，
并指出哪些组合含有扭度 2（twist-2）胶子 PDF 的信息。

接着讨论单圈修正。第 3 节分析规范链接的自能贡献
及其紫外和短距离行为的特定性质。
第 4 节以对前向和非前向情形都成立的形式给出胶子链接顶点修正的结果。
第 5 节讨论``盒图''。
由于此情形下的结果相当冗长，
我们只给出其中一部分，且仅限于前向情形。
第 6 节讨论胶子自能修正。

第 7 节的主题是胶子算符微扰演化及匹配条件的结构。
第 8 节总结全文。
